% TOP_PROBLEM: {"id": 3195, "title": "List Total Colouring Conjecture", "category_id": 3, "category": {"id": 3, "name": "graph_theory", "display_name": "Graph Theory", "description": "Problems involving graphs, networks, and their properties.", "slug": "graph-theory", "order_index": 3, "created_at": "2026-07-31T15:26:25.671Z"}, "status": "open", "problem_number": "OPG-56449", "set_id": 12}
% FIRST_POSTED: 2026-10-01
% ENTRY_KIND: statement_only
% UnsolvedMath ID 3195; source catalogue code OPG-56449
% !TeX program = lualatex
% Definition and sources only; this file is not an LLM solution attempt.
\documentclass[11pt,a4paper]{article}
\usepackage[margin=25mm]{geometry}
\usepackage{fontspec}
\setmainfont{lmroman10-regular.otf}[BoldFont=lmroman10-bold.otf,
 ItalicFont=lmroman10-italic.otf,BoldItalicFont=lmroman10-bolditalic.otf]
\usepackage{amsmath,amssymb,mathtools}
\usepackage{unicode-math}
\setmathfont{latinmodern-math.otf}
\AtBeginDocument{\let\setminus\smallsetminus}
\usepackage{xurl,hyperref}
\hypersetup{colorlinks=true,urlcolor=blue}
\setcounter{secnumdepth}{0}
\setlength{\parindent}{0pt}
\setlength{\parskip}{5pt}
\setlength{\emergencystretch}{3em}
\begin{document}
\section*{List Total Colouring Conjecture}\label{3195}
{\small UnsolvedMath ID 3195 · OPG-56449 · Graph Theory · OpenGarden}
\subsection{Definitions and mathematical statement}
Let $H$ be a finite loopless multigraph; parallel
edges are distinct elements. Its total graph
$T(H)$ is a simple graph on the disjoint union
$V(H)\sqcup E(H)$. Two distinct vertex-elements
are adjacent when joined by an edge of $H$;
two distinct edge-elements are adjacent when
their edges share an endpoint; a vertex-element
and an edge-element are adjacent when incident
in $H$.

A proper colouring of $T(H)$ therefore colours
both vertices and edges of $H$, separating
adjacent vertices, incident edges, and every
vertex from each incident edge. Let
$\chi(T(H))$ be the minimum common palette size.
For a simple graph $F$, a list assignment gives
each vertex $x$ a finite set $L(x)$ of allowed
colours. A proper list colouring chooses
$c(x)\in L(x)$ with different choices at
adjacent vertices. The list chromatic number
$\chi_\ell(F)$ is the least integer $q$ for
which every assignment with $|L(x)|\geq q$
admits such a colouring.

The List Total Colouring Conjecture states
\[
 \chi_\ell(T(H))=\chi(T(H))
 \qquad\text{for every finite loopless multigraph }H.
\]
Thus arbitrary separate lists for all vertices
and edges should suffice whenever their sizes
equal the ordinary total chromatic number.
The lists may overlap in any way and need not
come from one palette of that size. The inequality
$\chi_\ell\geq\chi$ is automatic from identical
lists; the reverse inequality is the target.
\subsection{Short English statement}
Do arbitrary colour lists of ordinary total-colouring size always suffice to colour every vertex and edge of a loopless multigraph?
\subsection{Sources}
\begin{itemize}
\item[S1] ULAM.AI and the UnsolvedMath Contributors, supplied problems.json, record 3195 (OPG-56449), OpenGarden collection. Catalogue identity and source transcription; CC BY 4.0.
\url{https://huggingface.co/datasets/ulamai/UnsolvedMath}.
\item[S2] Open Problem Garden, List Total Colouring Conjecture. Original problem and discussion; the mathematical formulation below is expanded from the supplied source record.
\url{http://www.openproblemgarden.org/op/list_total_colouring_conjecture}.
\end{itemize}
\end{document}
